% Ordered control priorities and the margin dynamics they act on.  Module roles,
% inputs, and integration boundaries belong to the overview subsection and are
% deliberately not restated here.
%
% Three claims this subsection must never make again; each contradicts the
% implementation or a later subsection.  (1) That margin "does not reset": every
% capture starts a fresh clock, and the coupling comes from serialized Draft
% execution.  (2) That pacing leaves a remainder to admission: no numeric state
% passes between the two modules.  (3) That either runtime test is an online form
% of \(\mu_i>0\): \(\mu_i\) is realized only at completion.  A scalar cost
% objective was also drafted and rejected; see AGENTS.md for why.

\subsection{Control Objective}
\label{sec:objective}

Capture Timeout prevention is the controller's hard goal, and it can pay two costs to meet it: a delayed next capture and a skipped optional Draft stage.
The delay is the more expensive by product policy, because it is what the user perceives directly and foreground usability takes priority.

Each capture starts a fresh Capture Timeout clock, but a single worker serializes the Drafts, so a capture's realized margin depends on when its predecessor's completed.
Capture \(i\) must complete its Draft path within the budget \(\Delta\) of its capture request; let \(\mu_i\) be its realized margin at that completion, \(\delta_i\) the interval between the capture requests of \(i-1\) and \(i\), and \(\pi_i\) the interval between the corresponding Draft-path completions.
These definitions give, for successive captures,
\begin{equation}
  \mu_i=\mu_{i-1}+\delta_i-\pi_i .
  \label{eq:margin-recursion}
\end{equation}
Only two quantities move the margin: admission acts directly on \(\pi_i\), pacing directly on \(\delta_{i+1}\).

A skipped stage suppresses a configured enhancement in the early Draft image, and final processing normally masks it.
Pacing therefore converts only a fixed fraction of the deficit it projects into a capture-availability delay, and admission independently applies its own live-budget test at each optional node.
The result is a safety-constrained, responsiveness-dominant trade-off, not the optimum of a cost model: the controller computes no quality utility and solves no global optimization.

Equation~\ref{eq:margin-recursion} is an accounting identity, not a control law: neither module tracks it, and each compares a decision-time forecast against the deadline window it can still read.
Neither \(\delta\) nor \(\pi\) is known when the decision that shapes it is taken, and every decision changes the backlog and thermal state that determine later ones.
Keeping \(\mu_i>0\) is therefore a design goal, not an invariant the controller can guarantee.
